\section{Mathematical Derivations and Proofs}
\label{app:derivations}

\subsection{Riemannian Geometry and Vector Field Derivations}
\label{app:riemannian-derivations}

The batch state lies on the product manifold
\begin{equation}
    \mathcal{M}
    =
    \left(\mathbb{S}^{d_S-1}\right)^B,
\end{equation}
equipped with the Frobenius inner product. Its tangent space at
$\mathbf{Z}$ is
\begin{equation}
    T_{\mathbf{Z}}\mathcal{M}
    =
    \left\{
        \mathbf{U}\in\mathbb{R}^{B\times d_S}
        :
        \mathbf{z}_i^\top\mathbf{u}_i=0
        \text{ for every }i
    \right\}.
    \label{eq:app-product-tangent}
\end{equation}
The row-wise projection in Eq.~\ref{eq:tangent-projection} follows directly:
\begin{equation}
    \Pi_{\mathbf{Z}}(\mathbf{U})_i
    =
    \mathbf{u}_i
    -
    (\mathbf{z}_i^\top\mathbf{u}_i)\mathbf{z}_i.
    \label{eq:app-batch-projection}
\end{equation}

\paragraph{Point-wise semantic gradient.}
For one normalized pair $(\mathbf{z},\boldsymbol{\tau})$, let
$c=\mathbf{z}^\top\boldsymbol{\tau}$ and
$\theta=\arccos(c)\in[0,\pi)$. We exclude the antipodal cut locus
$\theta=\pi$, where the shortest geodesic is not unique. The ambient
derivative of the squared geodesic potential is
\begin{equation}
    \nabla_{\mathbf{z}}
    \frac{1}{2}\theta^2
    =
    -\frac{\theta}{\sin\theta}
    \boldsymbol{\tau}.
\end{equation}
Projecting this derivative onto the tangent space gives
\begin{align}
    \operatorname{grad}_{\mathbb{S}}
    \frac{1}{2}\theta^2
    &=
    -\frac{\theta}{\sin\theta}
    \left(
        \boldsymbol{\tau}
        -
        c\mathbf{z}
    \right)
    \nonumber\\
    &=
    -\operatorname{Log}_{\mathbf{z}}(\boldsymbol{\tau}).
    \label{eq:app-semantic-single-gradient}
\end{align}
The continuous extension at $\theta=0$ is zero. Applying
Eq.~\ref{eq:app-semantic-single-gradient} to the batch mean in
Eq.~\ref{eq:semantic-potential} yields
\begin{equation}
    \operatorname{grad}_{\mathbb{S}}
    \mathcal{E}_{\mathrm{sem}}(\mathbf{Z},\mathbf{T})
    =
    -\frac{1}{B}
    \operatorname{Log}_{\mathbf{Z}}(\mathbf{T}).
    \label{eq:app-semantic-batch-gradient}
\end{equation}
The norm of each logarithmic-map row equals its geodesic distance because
\begin{equation}
    \left\lVert
        \boldsymbol{\tau}
        -
        c\mathbf{z}
    \right\rVert_2
    =
    \sin\theta,
    \qquad
    \left\lVert
        \operatorname{Log}_{\mathbf{z}}(\boldsymbol{\tau})
    \right\rVert_2
    =
    \theta.
    \label{eq:app-log-map-norm}
\end{equation}

\paragraph{Relational gradient.}
Let
\begin{equation}
    \mathbf{A}
    =
    \mathbf{Z}\mathbf{Z}^\top-\mathbf{G}_T.
\end{equation}
Both Gram matrices are symmetric, so $\mathbf{A}=\mathbf{A}^\top$. For a
perturbation $\mathbf{H}$, the differential of
Eq.~\ref{eq:geometric-potential} is
\begin{align}
    \mathrm{D}\mathcal{E}_{\mathrm{geo}}[\mathbf{H}]
    &=
    \frac{2}{B^2}
    \left\langle
        \mathbf{A},
        \mathbf{H}\mathbf{Z}^\top
        +
        \mathbf{Z}\mathbf{H}^\top
    \right\rangle_F
    \nonumber\\
    &=
    \frac{4}{B^2}
    \left\langle
        \mathbf{A}\mathbf{Z},
        \mathbf{H}
    \right\rangle_F.
    \label{eq:app-relational-differential}
\end{align}
Hence the ambient and Riemannian gradients are
\begin{align}
    \nabla_{\mathbf{Z}}\mathcal{E}_{\mathrm{geo}}
    &=
    \frac{4}{B^2}
    \mathbf{A}\mathbf{Z},
    \nonumber\\
    \operatorname{grad}_{\mathbb{S}}
    \mathcal{E}_{\mathrm{geo}}
    &=
    \frac{4}{B^2}
    \Pi_{\mathbf{Z}}(\mathbf{A}\mathbf{Z}).
    \label{eq:app-relational-gradient}
\end{align}

Combining Eqs.~\ref{eq:app-semantic-batch-gradient} and
\ref{eq:app-relational-gradient} gives
\begin{equation}
    \operatorname{grad}_{\mathbb{S}}\mathcal{E}
    =
    -\frac{\alpha}{B}
    \operatorname{Log}_{\mathbf{Z}}(\mathbf{T})
    +
    \frac{4\beta}{B^2}
    \Pi_{\mathbf{Z}}
    \!\left[
        (\mathbf{Z}\mathbf{Z}^\top-\mathbf{G}_T)\mathbf{Z}
    \right].
    \label{eq:app-total-riemannian-gradient}
\end{equation}
Multiplying by $-B$ recovers the vector direction in
Eq.~\ref{eq:semantic-vector-direction}. This scaling removes the $1/B$ from
the batch mean without changing the descent direction.

\subsection{Finite-Horizon Flow Derivation}
\label{app:finite-horizon-derivation}

Consider first the teacher-conditioned gradient flow in an auxiliary time
$\tau$,
\begin{equation}
    \frac{d\mathbf{Z}}{d\tau}
    =
    \mathbf{V}(\mathbf{Z},\mathbf{T},\mathbf{G}_T)
    =
    -B\operatorname{grad}_{\mathbb{S}}\mathcal{E}.
    \label{eq:app-auxiliary-flow}
\end{equation}
This flow generally reaches a stationary point only as
$\tau\rightarrow\infty$. To fit that horizon into student depth
$t\in[0,1]$, define
\begin{equation}
    \tau(t)
    =
    \int_0^t\frac{s(u)}{R(u)}\,du,
    \qquad
    R(t)=\int_t^1s(u)\,du.
    \label{eq:app-time-change}
\end{equation}
Since $\dot{R}(t)=-s(t)$,
\begin{equation}
    \frac{d}{dt}\log R(t)
    =
    -\frac{s(t)}{R(t)},
\end{equation}
and therefore
\begin{equation}
    \tau(t)
    =
    \log\frac{R(0)}{R(t)}.
    \label{eq:app-warped-time}
\end{equation}
For any schedule with $R(t)>0$ for $t<1$ and $R(1)=0$,
$\tau(t)\rightarrow\infty$ as $t\rightarrow1$. Applying the chain rule to
Eq.~\ref{eq:app-auxiliary-flow} gives
\begin{equation}
    \frac{d\mathbf{Z}}{dt}
    =
    \frac{d\tau}{dt}
    \frac{d\mathbf{Z}}{d\tau}
    =
    \frac{s(t)}{R(t)}
    \mathbf{V}(\mathbf{Z},\mathbf{T},\mathbf{G}_T),
\end{equation}
which is Eq.~\ref{eq:finite-horizon-flow}.

\paragraph{Schedules.}
For $s(t)=t^p$ with $p>-1$,
\begin{equation}
    R(t)
    =
    \frac{1-t^{p+1}}{p+1}.
    \label{eq:app-power-remaining-mass}
\end{equation}
The constant schedule is $p=0$, while the default linear schedule is $p=1$
and gives $R(t)=(1-t^2)/2$.

\paragraph{Discrete step fraction.}
The warped time elapsed between two student depths is
\begin{align}
    \Delta\tau_\ell
    &=
    \int_{t_\ell}^{t_{\ell+1}}
    \frac{s(u)}{R(u)}\,du
    \nonumber\\
    &=
    \log\frac{R(t_\ell)}{R(t_{\ell+1})}.
    \label{eq:app-warped-layer-interval}
\end{align}
Under the instance-only auxiliary flow, the distance remaining after this
interval is multiplied by $\exp(-\Delta\tau_\ell)$. The fraction covered is
therefore
\begin{align}
    \rho_\ell
    &=
    1-\exp(-\Delta\tau_\ell)
    \nonumber\\
    &=
    1-\frac{R(t_{\ell+1})}{R(t_\ell)},
    \label{eq:app-step-fraction}
\end{align}
which recovers Eq.~\ref{eq:finite-step-fraction}. For the default schedule and
$t_\ell=\ell/L$,
\begin{equation}
    \rho_\ell
    =
    \frac{2\ell+1}{L^2-\ell^2},
    \qquad
    \rho_{L-1}=1.
    \label{eq:app-linear-step-fraction}
\end{equation}
For the point-wise field, this fraction reproduces the exact geodesic progress
over the interval. For the coupled semantic--relational field, GeoODE-KD uses
the same $\rho_\ell$ as a local exponential-map step; it is a discretization,
not an exact integration of the changing nonlinear field.

The exponential map also preserves the discrete target norm. If
$\mathbf{v}\in T_{\mathbf{z}}\mathbb{S}^{d_S-1}$, then
$\mathbf{z}^\top\mathbf{v}=0$, and Eq.~\ref{eq:sphere-exp-map} gives
\begin{equation}
    \left\lVert
        \operatorname{Exp}_{\mathbf{z}}(\mathbf{v})
    \right\rVert_2^2
    =
    \cos^2\lVert\mathbf{v}\rVert_2
    +
    \sin^2\lVert\mathbf{v}\rVert_2
    =1.
    \label{eq:app-exp-map-norm}
\end{equation}

\subsection{Proofs of the Flow Properties}
\label{app:proofs}

The following statements assume $s(t)\geq0$, $R(t)>0$ for $t<1$, and that no
student--teacher pair lies at the antipodal cut locus. The logarithmic map is
defined by continuous extension when a student state already equals its target.

\paragraph{Proposition 1 (norm preservation).}
Every row of a solution to Eq.~\ref{eq:finite-horizon-flow} remains on the unit
hypersphere.

\paragraph{Proof.}
Both terms in Eq.~\ref{eq:semantic-vector-direction} are tangent: the log map
lies in $T_{\mathbf{z}_i}\mathbb{S}^{d_S-1}$, and the relational term is
explicitly projected by $\Pi_{\mathbf{Z}}$. Thus
$\mathbf{z}_i^\top\dot{\mathbf{z}}_i=0$ and
\begin{equation}
    \frac{d}{dt}
    \lVert\mathbf{z}_i(t)\rVert_2^2
    =
    2\mathbf{z}_i(t)^\top\dot{\mathbf{z}}_i(t)
    =0.
\end{equation}
Since $\lVert\mathbf{z}_i(0)\rVert_2=1$, its norm remains one. \hfill$\square$

\paragraph{Proposition 2 (energy descent).}
Along the ideal continuous flow, the teacher-conditioned potential is
non-increasing for every $t<1$.

\paragraph{Proof.}
By Eq.~\ref{eq:app-total-riemannian-gradient},
$\mathbf{V}=-B\operatorname{grad}_{\mathbb{S}}\mathcal{E}$. The Riemannian
chain rule gives
\begin{align}
    \frac{d\mathcal{E}}{dt}
    &=
    \left\langle
        \operatorname{grad}_{\mathbb{S}}\mathcal{E},
        \dot{\mathbf{Z}}
    \right\rangle_F
    \nonumber\\
    &=
    -\frac{B\,s(t)}{R(t)}
    \left\lVert
        \operatorname{grad}_{\mathbb{S}}\mathcal{E}
    \right\rVert_F^2
    \leq0.
\end{align}
This proves Eq.~\ref{eq:energy-descent}. \hfill$\square$

\paragraph{Corollary 1 (point-wise geodesic contraction).}
For $\alpha=1$ and $\beta=0$, the geodesic distance
$\theta(t)=d_g(\mathbf{z}(t),\boldsymbol{\tau})$ satisfies
\begin{equation}
    \theta(t)
    =
    \theta(0)\frac{R(t)}{R(0)},
    \qquad
    \lim_{t\rightarrow1}\theta(t)=0.
    \label{eq:app-pointwise-contraction}
\end{equation}

\paragraph{Proof.}
For one pair, define $e(t)=\tfrac{1}{2}\theta(t)^2$. From
Eq.~\ref{eq:app-semantic-single-gradient} and the point-wise form of
Eq.~\ref{eq:finite-horizon-flow},
\begin{align}
    \dot{e}(t)
    &=
    \left\langle
        -\operatorname{Log}_{\mathbf{z}}(\boldsymbol{\tau}),
        \frac{s(t)}{R(t)}
        \operatorname{Log}_{\mathbf{z}}(\boldsymbol{\tau})
    \right\rangle
    \nonumber\\
    &=
    -\frac{s(t)}{R(t)}\theta(t)^2,
\end{align}
where Eq.~\ref{eq:app-log-map-norm} is used in the last line. For
$\theta(t)>0$, $\dot{e}=\theta\dot{\theta}$, so
\begin{equation}
    \frac{d}{dt}\log\theta(t)
    =
    -\frac{s(t)}{R(t)}
    =
    \frac{d}{dt}\log R(t).
\end{equation}
Integrating from $0$ to $t$ proves Eq.~\ref{eq:app-pointwise-contraction}. If
$\theta$ reaches zero earlier, the log-map vector is zero and the state remains
at the target. \hfill$\square$

If the observed student trajectory begins at the first Transformer layer
$t_1=1/L$ rather than at $t=0$, the same argument starts from $t_1$ and gives
\begin{equation}
    \theta(t)
    =
    \theta(t_1)\frac{R(t)}{R(t_1)},
    \qquad t\geq t_1.
    \label{eq:app-first-layer-contraction}
\end{equation}

Since $\theta(t)\in[0,\pi)$ and $\dot{\theta}(t)\leq0$, teacher cosine is also
non-decreasing:
\begin{equation}
    \frac{d}{dt}\cos\theta(t)
    =
    \frac{s(t)}{R(t)}
    \theta(t)\sin\theta(t)
    \geq0.
    \label{eq:app-cosine-monotonicity}
\end{equation}

\paragraph{Corollary 2 (exact sampled profile).}
For $0\leq a\leq1$, the curve
\begin{equation}
    \gamma(a)
    =
    \operatorname{Exp}_{\mathbf{z}}
    \!\left(
        a\operatorname{Log}_{\mathbf{z}}(\boldsymbol{\tau})
    \right)
\end{equation}
is the shortest geodesic from $\mathbf{z}$ to $\boldsymbol{\tau}$ and satisfies
\begin{equation}
    d_g(\gamma(a),\boldsymbol{\tau})
    =
    (1-a)d_g(\mathbf{z},\boldsymbol{\tau}).
    \label{eq:app-slerp-distance}
\end{equation}
Setting $a=\rho_\ell$ in the instance-only target gives
\begin{equation}
    \theta_{\ell+1}
    =
    (1-\rho_\ell)\theta_\ell
    =
    \frac{R(t_{\ell+1})}{R(t_\ell)}\theta_\ell.
    \label{eq:app-discrete-contraction}
\end{equation}
The factors telescope, so the predicted states reproduce
Eq.~\ref{eq:app-pointwise-contraction} at every sampled depth and the last
target equals $\boldsymbol{\tau}$ because $\rho_{L-1}=1$.

\paragraph{Scope of the guarantees.}
Norm preservation and energy descent describe the ideal continuous field.
Exact arrival and the sampled profile require the instance-only setting
$\alpha=1,\beta=0$. With $\beta>0$, the native teacher Gram target can compete
with the projected point-wise target, so individual distances need not decrease
at every depth. Finally, the realized Transformer trajectory satisfies the
sampled profile only to the extent that it minimizes
$\mathcal{L}_{\mathrm{evel}}$; the propositions do not assume zero
discretization or optimization error.
